\chapter{符号、术语与主要定理依赖表}
\label{app:notation-dependencies}
\label{chap:appendix-notation-dependencies}

\begin{chapterguide}
本附录统一前三编已经使用的记号、同名定理版本和默认假设。全书后续卷涉及的拓扑、测度、函数空间与偏微分方程记号保留既定栏目，待相应正文定稿后逐项补入，避免符号表先于定义冻结。
\end{chapterguide}

\section{数集、拓扑、测度和概率记号}

\begin{description}[style=nextline,leftmargin=4.2em]
 \item[\(\N,\Npos,\Z,\Q,\R,\C\)]
 自然数、正整数、整数、有理数、实数、复数。本书固定 \(\N=\{0,1,2,\ldots\}\)，并以 \(\Npos=\{1,2,3,\ldots\}\) 表示正整数集。
 \item[\(\R_{>0},\R_{\ge0}\)]
 正实数与非负实数；其他数集使用同样下标。
 \item[\(\overline{\R}\)]
 扩充实数 \(\R\cup\{-\infty,+\infty\}\)。\(\pm\infty\) 是序端点，不是实数；未定式不按普通代数运算处理。
 \item[{\((a,b),[a,b]\)}]
 开区间与闭区间；半开区间写为 \((a,b]\)、\([a,b)\)。端点可为扩充实数时，包含无穷端点的方括号只在序拓扑邻域中使用。
 \item[\(U(a,\delta),U^\circ(a,\delta)\)]
 普通邻域与删心邻域：
 \[
  U(a,\delta)=\{x:\abs{x-a}<\delta\},\quad
  U^\circ(a,\delta)=\{x:0<\abs{x-a}<\delta\}.
 \]
 \item[\(\abs{x},\norm{x}\)]
 前者为实数绝对值；后者为向量或函数空间范数，必须由上下文或首次出现处指定。
 \item[\(\sup A,\inf A,\max A,\min A\)]
 上确界、下确界、最大元、最小元。前两者不必属于 \(A\)，后两者必须属于 \(A\)。
 \item[\(\lfloor x\rfloor,\lceil x\rceil,\{x\}\)]
 下取整、上取整和小数部分；\(\{x\}=x-\lfloor x\rfloor\)。
 \item[\(\mathcal P(X),\lvert X\rvert\)]
 幂集与基数。本书不把“可数”单独用作正式分类词：与 \(\N\) 等势称“可数无限”，有限或可数无限称“至多可数”。
 \item[\(f|_E,f^{-1}(A)\)]
 限制映射与集合原像。\(f^{-1}\) 只有在上下文明确为双射逆映射时才表示逆函数。
\end{description}

数列写作 \((a_n)\) 或 \((a_n)_{n\ge1}\)，子列写作 \((a_{n_k})\)，并默认
\[
 n_1<n_2<\cdots.
\]
部分极限、上极限和下极限分别写作“子列极限”、
\(\limsup a_n\)、\(\liminf a_n\)。函数极限若定义域限制重要，写作
\[
 \lim_{\substack{x\to a\\x\in D}}f(x).
\]

\subsection*{拓扑、流形、测度和概率记号}

\begin{remark}[拓扑记号]
下列记号分别表示闭包、内部、边界、距离和开球：
\[
 \overline A,\quad A^\circ,\quad\partial A,\quad
 d(x,y),\quad B(x,r),
\]
它们在欧氏空间、度量空间和一般拓扑空间各章逐层定义。前三编只在已经明确给出含义的实直线情形使用，不把本表当作定义的替代。
\end{remark}

拓扑基写作 \(\mathcal B\)，生成的拓扑写作 \(\tau\)；网以有向集为指标，滤子与超滤子分别写作 \(\mathcal F,\mathcal U\)。流形的切空间写作 \(T_pM\)，切丛写作 \(TM\)。Borel \(\sigma\)-代数写作 \(\mathcal B(X)\)，测度空间写作 \((X,\mathcal A,\mu)\)。概率空间写作 \((\Omega,\mathcal F,\mathbb P)\)，期望与条件期望写作 \(\mathbb E X\) 和 \(\mathbb E[X\mid\mathcal G]\)。这些约定只有在相应正文给出定义后才可用于推理。

\section{函数空间、算子和偏微分方程记号}

\begin{remark}[定量连续性记号]
\[
 \operatorname{Lip}(f),\qquad
 \omega_f(t),\qquad C^{0,\alpha}
\]
分别用于 Lipschitz 常数、连续模和 Hölder 类。前三编只在实值函数与实数子集上使用；后续推广到度量值域时改用距离。
\end{remark}

\begin{remark}[函数空间与弱导数记号]
常用函数空间记作
\[
 C^k,\quad C_c^\infty,\quad L^p,\quad \ell^p,\quad
 W^{k,p},\quad H^s.
\]
等号在 \(L^p\) 与 Sobolev 空间中表示几乎处处等价类。连续对偶写作 \(X^*\)，有界线性算子空间写作 \(\mathcal L(X,Y)\)，算子范数写作 \(\|T\|\)，谱写作 \(\sigma(T)\)。分布导数使用 \(D^\alpha u\)，弱解所在空间和测试函数空间必须在具体方程旁给出；本表不设任何默认边界条件。
\end{remark}

\section{同名定理的不同版本}

\begin{description}[style=nextline,leftmargin=4em]
 \item[Heine 定理]
 本书指函数极限的序列刻画。输入点列必须属于定义域，有限考察点处还要排除恒取中心的项。
 \item[Heine--Borel 定理]
 在实线中，闭且有界等价于每个开覆盖有有限子覆盖。高维欧氏空间有对应版本；一般度量空间中“闭且有界”不再充分。
 \item[Heine--Cantor 定理]
 紧集上的连续函数一致连续。值域可推广到度量空间。
 \item[Bolzano--Weierstrass 定理]
 前三编指有界实数列有收敛子列；等价的集合版本是无限有界实数集有聚点。
 \item[Cauchy 准则]
 数列版本比较尾部两项；函数极限版本比较同一删心邻域中两点；二者的充分性都依赖值域完备。
 \item[介值定理与零点定理]
 零点定理是端点异号情形；介值定理对 \(f-y\) 应用零点定理得到。连续像保持连通性是其结构版本。
 \item[单调收敛定理]
 前三编指单调有界实数列收敛。测度论中的“单调收敛定理”是积分定理，后续引用必须加上“数列”或“积分”限定。
\end{description}

\section{主要定理依赖图}

数列极限的定义和代数性质只使用实数的序与域运算；完备性负责保证某些候选极限确实存在。因此依赖关系不是一条把所有内容串起的单链。前三编的主要分支为
\[
\begin{array}{c}
\text{逻辑、集合与有理数}\\[2mm]
\downarrow\\[-1mm]
\text{完备有序域的公理描述}\quad\longleftarrow\quad
\begin{matrix}
\text{Dedekind 分割构造}\\
\text{Cauchy 列商构造}
\end{matrix}\\[3mm]
\downarrow\\[-1mm]
\text{Archimedes 性质、整数部分、有理稠密}\\[2mm]
\swarrow\qquad\downarrow\qquad\searrow\\[-1mm]
\begin{matrix}\text{单调有界}\\\text{收敛}\end{matrix}
\quad
\begin{matrix}\text{区间套与}\\\text{Bolzano--Weierstrass}\end{matrix}
\quad
\begin{matrix}\text{Cauchy}\\\text{完备性}\end{matrix}.
\end{array}
\]
另一方面，
\[
\text{数列极限定义与运算}
\longrightarrow
\text{子列、上下极限、Cauchy 判据}
\longrightarrow
\text{函数极限与连续性}.
\]
两个分支在“存在性”处相交：单调数列、部分极限和 Cauchy 列的收敛都调用相应的完备性形式。随后还有以下侧向依赖：
\begin{itemize}
 \item 函数极限的 Cauchy 准则在充分性方向使用值域完备性；
 \item 单调函数的左右极限使用确界性质，其间断点至多可数还使用有理数稠密性；
 \item 闭区间有限覆盖可由确界、区间套或序列紧致得到；
 \item 极值与介值分别来自紧致性和连通性；
 \item Heine--Cantor 由紧致性与连续性推出，一致延拓还需要值域完备。
\end{itemize}
依赖图只记录可以引用的先后关系，不代替正文证明。若某个证明沿箭头反向调用后文结论，就必须另给不循环的证明。

\section{默认假设、例外和选择原则}

\begin{enumerate}
 \item 未注明时，数列为实数列，函数为实变量实值函数；推广到度量值域时显式写出空间。
 \item 函数极限 \(\lim_{x\to a}f(x)\) 默认 \(a\) 是定义域聚点；连续性要求 \(a\in D\)。
 \item “区间”允许开、闭、半开、无界与单点退化情形；涉及左右极限时另行要求相应单侧聚点。
 \item “紧集”在前三编指实线中具有有限覆盖性质的集合，等价于闭且有界；后续拓扑卷以开覆盖定义为准。
 \item “完备”必须说明相对于哪个距离或一致结构；\(\Q\) 在通常距离下不完备。
 \item 上确界、下确界的正文定理默认集合非空并有相应界；空集与扩充实值约定须明确说明。
 \item 常数 \(C,c,M\) 可逐行改变时写明“存在与主要变量无关的常数”；依赖关系用 \(C(\alpha,p,\ldots)\) 标注。
 \item 从固定枚举中取第一个合格元素属于规范选择；从任意非空集合族同时取元可能使用选择公理。所用强度按附录 A 的分类说明。
 \item 本表中出现的后卷符号只登记字体和含义，不构成定义，也不得作为前卷已证明结果引用。
\end{enumerate}

\section{编写阶段登记表的整理规则}

本附录是给读者查阅的成书内容，不承担工程日志职能。编写阶段的活符号表、术语表和定理依赖表分别维护在
\texttt{docs/NOTATION.md}、\texttt{docs/TERMINOLOGY.md} 与
\texttt{docs/DEPENDENCIES.md} 中；它们可记录尚未进入正文的候选项、文件位置和审校状态。

整理进本附录时遵守三条规则：
\begin{enumerate}
 \item 只收入已经在正文中定义，或确有必要提前统一字体的记号；后者必须明说不能替代定义。
 \item 同一名称有多个版本时，列出对象范围和假设差别，不以“同理”掩盖维数、完备性或可测性条件。
 \item 依赖箭头以正文中实际可引用的证明为准。历史报告和计划中的拟定顺序不能覆盖现行数学依赖。
\end{enumerate}
因此，工程登记表可以频繁更新，而本附录只在有关理论定稿后整理；二者职责不同。

\section*{习题}
\addcontentsline{toc}{section}{习题}

\begin{enumerate}[label=\textbf{E.\arabic*.}]
 \exerciseitem{E.1} 区分 \(\sup A\) 与 \(\max A\)，并为四种“均存在且相等、只有确界、两者都不存在、下确界对应情形”各给一例。
 \exerciseitem{E.2} 写出数列极限、函数删心极限、点连续和一致连续的量词形式，比较变量依赖。
 \exerciseitem{E.3} 解释 \(f^{-1}\) 作为原像与逆函数的两种含义，并给出不会混淆的上下文写法。
 \exerciseitem{E.4} 画出从确界原理到 Cauchy 完备、再到函数 Cauchy 准则的依赖链。
 \exerciseitem{E.5} 比较 Heine、Heine--Borel 与 Heine--Cantor 三个名称对应的结论和假设。
 \exerciseitem{E.6} 说明“闭且有界”“紧致”“序列紧致”在前三编中的关系，以及为何不得无条件推广到后文空间。
 \projectexerciseitem{E.7} \ 【前置：前三编；预计用时：60分钟】为\chapterref{chap:function-limits}至\chapterref{chap:uniform-continuity}制作一页符号与依赖速查表，要求每个定理列出输入对象、关键假设和输出结论。
 \openexerciseitem{E.8} 选择一个后续预留栏目，提出定稿时必须统一的五项记号决策；不得自行把决定写入全书规范。
\end{enumerate}
